\documentclass[11pt]{article}

\usepackage[margin=1in]{geometry}
\usepackage{fontspec}
\usepackage{microtype}
\usepackage{xcolor}
\usepackage{titlesec}
\usepackage{enumitem}
\usepackage{tabularx}
\usepackage{array}
\usepackage{booktabs}
\usepackage{fancyhdr}
\usepackage{hyperref}
\usepackage{longtable}

\setmainfont{TeX Gyre Pagella}
\setsansfont{TeX Gyre Heros}
\hypersetup{colorlinks=true, linkcolor=blue!50!black, urlcolor=blue!50!black}
\definecolor{purduegold}{HTML}{CEB888}
\definecolor{navy}{HTML}{101050}
\definecolor{softgray}{HTML}{F4F4F4}
\titleformat{\section}{\Large\bfseries\sffamily\color{navy}}{\thesection}{0.6em}{}
\titleformat{\subsection}{\large\bfseries\sffamily\color{navy}}{\thesubsection}{0.6em}{}
\titleformat{\subsubsection}{\bfseries\sffamily\color{navy}}{\thesubsubsection}{0.6em}{}
\setlist[itemize]{topsep=4pt,itemsep=2pt,parsep=1pt,leftmargin=1.25em}
\setlist[enumerate]{topsep=4pt,itemsep=2pt,parsep=1pt,leftmargin=1.5em}
\newcolumntype{Y}{>{\raggedright\arraybackslash}X}
\newcommand{\keybox}[1]{\vspace{0.35em}\noindent\fcolorbox{navy}{softgray}{\parbox{0.96\linewidth}{#1}}\vspace{0.35em}}
\sloppy

\pagestyle{fancy}
\setlength{\headheight}{14pt}
\fancyhf{}
\lhead{EDCI 59100-016}
\rhead{Week 6 Lecture Notes}
\cfoot{\thepage}
\renewcommand{\headrulewidth}{0.4pt}

\begin{document}

\begin{titlepage}
\centering
\vspace*{1.2in}
{\Large EDCI 59100-016: Educational Research and Literature Review\par}
\vspace{0.6in}
{\Huge\bfseries Week 6 Comprehensive Lecture Notes\par}
\vspace{0.25in}
{\LARGE Revising Toward a Submission-Ready Literature Review Manuscript\par}
\vfill
{\Large Summer 2026\par}
\vfill
\begin{minipage}{0.78\linewidth}
\textbf{Major themes covered:} manuscript coherence, methodology alignment, scholarly voice, clarity, precision, interpretation, implications, limitations, next steps, reviewer expectations, peer review, publication fit, revision workflow, and final manuscript rubric expectations.
\end{minipage}
\vfill
Prepared from the saved Week 6 course materials and final manuscript assignment/rubric.
\end{titlepage}

\tableofcontents
\newpage

\section{Week 6 Overview}

Week 6 shifts the course from developing separate parts of a literature review to revising the manuscript as a whole. The central question is no longer only whether sections exist. The central question is whether the manuscript works as a coherent, methodologically aligned, scholarly argument.

\keybox{\textbf{Central task for Week 6:} revise the literature review manuscript so that purpose, organization, analysis, voice, interpretation, limitations, and next steps fit together as one readable scholarly document.}

\subsection{Learning Objectives}

By the end of this week, students should be able to:

\begin{enumerate}
  \item evaluate whether manuscript sections connect into a coherent whole;
  \item identify places where organization, transitions, or repeated ideas weaken the throughline;
  \item revise claims for clarity, precision, and scholarly tone;
  \item check whether the manuscript reflects the selected literature review methodology;
  \item develop interpretation, implications, limitations, and next steps in a way that fits the current stage of the review;
  \item use peer review and a reviewer lens to identify high-impact revisions;
  \item compare the manuscript with expectations for journals, conferences, books, or other venues;
  \item prepare the final manuscript around the course rubric.
\end{enumerate}

\subsection{The Week 6 Revision Mindset}

Revision is not simply proofreading. Proofreading checks grammar, spelling, citation details, and surface errors. Week 6 revision asks deeper questions:

\begin{itemize}
  \item Does the introduction establish a clear direction?
  \item Does the methodology or approach explain how the literature is being reviewed?
  \item Do body sections synthesize sources rather than list them?
  \item Do headings and transitions guide the reader through the argument?
  \item Does the manuscript show why the literature matters?
  \item Are limitations and next steps realistic and methodologically aligned?
\end{itemize}

\section{From Sections to a Whole Manuscript}

\subsection{What Changes in Week 6}

Earlier work in the course emphasized finding sources, evaluating literature, selecting a review type, developing synthesis, and drafting sections. Week 6 assumes those pieces are now available in some form. The work becomes integrative: the writer examines how the pieces operate together.

\begin{center}
\begin{tabularx}{0.95\linewidth}{>{\raggedright\arraybackslash}p{0.28\linewidth}Y}
\toprule
\textbf{Earlier drafting question} & \textbf{Week 6 revision question} \\
\midrule
Do I have sources? & Are sources used to support a coherent synthesis? \\
Do I have sections? & Do sections build on one another logically? \\
Did I summarize studies? & Did I explain patterns, tensions, relationships, or gaps? \\
Did I name my method? & Does the manuscript actually behave like that review type? \\
Did I write a conclusion? & Do I explain interpretation, contribution, limitations, and next steps? \\
\bottomrule
\end{tabularx}
\end{center}

\subsection{The Manuscript as an Argument}

A literature review manuscript should do more than gather relevant studies. It should make a scholarly contribution by organizing the literature in a purposeful way. The contribution may be a thematic interpretation, an integrated framework, a conceptual clarification, a map of a field, a structured answer to a review question, or another methodologically appropriate product.

The manuscript therefore needs a visible line of reasoning. Readers should be able to answer:

\begin{itemize}
  \item What is the review about?
  \item Why does this topic matter?
  \item What kind of review is being conducted?
  \item How is the literature organized?
  \item What patterns or insights emerge?
  \item What remains limited, incomplete, or still developing?
\end{itemize}

\section{Coherence Across the Manuscript}

\subsection{What Coherence Means}

Coherence refers to the logical and consistent connection among parts of a manuscript. In writing, coherence means the reader can follow how ideas connect across sentences, paragraphs, sections, and the whole paper.

A coherent literature review:

\begin{itemize}
  \item communicates a clear purpose;
  \item connects ideas across sections;
  \item organizes themes logically;
  \item guides the reader through the author's thinking;
  \item makes the contribution visible;
  \item avoids disconnected summary.
\end{itemize}

Coherence is related to cohesion, but the two are not identical. Cohesion often involves sentence-level links such as pronouns, transition words, and repeated key terms. Coherence is larger: the whole manuscript fits together.

\subsection{Weak and Strong Coherence}

A weak transition only announces the next section:

\begin{quote}
This section focuses on peer feedback. The next section focuses on writing groups.
\end{quote}

A stronger transition explains the relationship:

\begin{quote}
Building on the challenges described in the previous section, this section examines peer feedback as a support structure that helps make academic expectations more visible.
\end{quote}

The second version tells the reader why the next section belongs where it is. This is the kind of revision Week 6 emphasizes.

\section{A Coherence Revision Process}

\subsection{Step 1: Revisit the Purpose}

Write the manuscript's central purpose in one sentence. Then ask whether each section clearly contributes to that purpose. If a section does not support the purpose, it may need to be revised, moved, combined, or removed.

\keybox{\textbf{Revision prompt.} My review examines \underline{\hspace{1.2in}} in order to show \underline{\hspace{1.2in}}. The manuscript contributes by \underline{\hspace{2in}}.}

\subsection{Step 2: Create a Reverse Outline}

A reverse outline describes what the existing manuscript actually does. It is different from a planning outline because it is based on the current draft.

\begin{enumerate}
  \item List each section or major paragraph.
  \item Write one sentence describing its function, not just its topic.
  \item Mark whether the section introduces, defines, compares, maps, synthesizes, interprets, or concludes.
  \item Identify gaps, repetition, abrupt shifts, and weak connections.
\end{enumerate}

\subsection{Step 3: Check Section Order}

Section order should help the reader move through the review. Ask:

\begin{itemize}
  \item Does the manuscript begin with context before moving to specialized analysis?
  \item Are definitions or theoretical concepts introduced before they are applied?
  \item Are themes sequenced from broad to specific, chronological to current, problem to response, or another logical pattern?
  \item Does each section prepare the reader for the next one?
\end{itemize}

\subsection{Step 4: Strengthen Transitions}

Transitions should explain relationships. Common transition purposes include:

\begin{center}
\begin{tabularx}{0.95\linewidth}{>{\raggedright\arraybackslash}p{0.25\linewidth}Y}
\toprule
\textbf{Transition purpose} & \textbf{What the transition should show} \\
\midrule
Addition & The next idea extends or deepens the previous one. \\
Contrast & The next idea complicates, challenges, or qualifies the previous one. \\
Cause or effect & The next idea explains why something happens or what follows from it. \\
Methodological shift & The next section changes the lens, evidence base, or analytic function. \\
Synthesis & Multiple sources or themes are being connected into a broader claim. \\
\bottomrule
\end{tabularx}
\end{center}

\subsection{Step 5: Identify Repetition}

Repetition is not always bad. Key terms and central claims should reappear. The problem is repetition that does not develop the argument. If a point repeats without adding evidence, interpretation, or connection, revise it.

\subsection{Step 6: Check the Throughline}

The throughline is the main idea running across the manuscript. A reader should be able to summarize the manuscript in two or three sentences. If the reader would only remember a list of disconnected topics, the throughline needs strengthening.

\section{Scholarly Voice, Clarity, and Precision}

\subsection{What Scholarly Voice Does}

Scholarly voice communicates ideas clearly, carefully, and professionally. It emphasizes evidence and interpretation rather than personal preference or unsupported assertion.

In a literature review, scholarly voice should:

\begin{itemize}
  \item keep the focus on ideas and patterns rather than only individual authors;
  \item support claims with evidence;
  \item integrate sources smoothly;
  \item use cautious and accurate language;
  \item avoid vague terms and overstatement.
\end{itemize}

\subsection{Moving from Author-by-Author Summary to Synthesis}

Less developed:

\begin{quote}
Smith found that feedback helps students. Lee also found that feedback helps students.
\end{quote}

More developed:

\begin{quote}
Research consistently shows that feedback supports student learning by making expectations more visible and actionable.
\end{quote}

The stronger version emphasizes the pattern across the literature. Specific citations would still be included in the manuscript, but the sentence is organized around the idea rather than around a list of authors.

\subsection{Strengthening Claims}

A strong claim is specific and meaningful. It tells the reader what is happening and why it matters.

\begin{center}
\begin{tabularx}{0.95\linewidth}{>{\raggedright\arraybackslash}p{0.32\linewidth}Y}
\toprule
\textbf{Less precise claim} & \textbf{Stronger claim} \\
\midrule
Feedback is important in learning. & Feedback improves student engagement by making expectations visible and actionable. \\
Many studies discuss teacher preparation. & Recent teacher preparation research emphasizes clinical practice, mentor feedback, and structured reflection as linked supports. \\
Technology affects education. & Digital tools can expand access, but their effect depends on instructional design, support, and institutional context. \\
\bottomrule
\end{tabularx}
\end{center}

\subsection{Improving Precision}

Precision means choosing words that match the exact meaning intended. Replace vague words such as \textit{things}, \textit{good}, \textit{bad}, \textit{a lot}, \textit{important}, and \textit{helpful} with more specific terms.

\keybox{\textbf{Revision prompt.} Find one vague word in the manuscript. Replace it with a term that names the specific concept, outcome, relationship, population, or process.}

\subsection{Deepening Analysis}

Analysis explains what evidence suggests. It is not enough to state that studies found something. A literature review should show what those findings mean together.

\begin{quote}
Feedback increases engagement, suggesting that making expectations visible plays a key role in supporting learning.
\end{quote}

The phrase after the comma performs analysis because it interprets the significance of the finding.

\subsection{Refining Tone}

Scholarly tone should be confident but careful. Avoid claiming that evidence \textit{proves} something unless the evidence truly supports that level of certainty. Use language such as \textit{suggests}, \textit{indicates}, \textit{appears to}, \textit{may}, or \textit{is associated with} when the evidence is conditional.

\section{Methodology Alignment}

\subsection{What Alignment Means}

Methodology alignment means the purpose, structure, and analysis of the manuscript match the selected review type. The manuscript should not merely name a review type; it should operate according to that type.

\keybox{\textbf{Diagnostic question:} Does my manuscript reflect the type of work my selected review methodology requires?}

\subsection{Alignment Across Review Types}

\begin{center}
\begin{tabularx}{0.98\linewidth}{>{\raggedright\arraybackslash}p{0.2\linewidth}Y Y}
\toprule
\textbf{Review type} & \textbf{Main purpose} & \textbf{What the manuscript should show} \\
\midrule
Narrative & Interpret a body of literature around a topic or problem. & Thematic organization, scholarly storyline, explanation of patterns and meanings. \\
Integrative & Bring together diverse empirical, conceptual, or theoretical sources. & Connections across types of evidence, tensions, integrated themes, or a framework. \\
Theoretical & Examine, compare, extend, or critique theories. & Definitions of theories, comparison of assumptions, development of a theoretical argument. \\
Conceptual & Clarify a concept and how it is understood. & Definitions, dimensions, categories, relationships, or conceptual structure. \\
Scoping & Map the field and identify what research exists. & Categories, distributions, gaps, evidence clusters, and boundaries of the field. \\
Systematic & Answer a focused question through transparent procedures. & Search criteria, selection process, structured findings, and methodologically organized results. \\
\bottomrule
\end{tabularx}
\end{center}

\subsection{Signs of Misalignment}

Misalignment may appear when:

\begin{itemize}
  \item the manuscript claims to be a scoping review but mainly offers personal interpretation without mapping the field;
  \item the manuscript claims to be theoretical but does not define or compare theories;
  \item the manuscript claims to be integrative but keeps empirical and conceptual sources separate;
  \item the manuscript claims to be systematic but does not explain selection procedures;
  \item the organization does not match the stated purpose.
\end{itemize}

\subsection{How to Revise Misalignment}

If the manuscript is too descriptive, add analysis that explains patterns, relationships, or meaning across studies. If it is too interpretive for the selected method, revise toward mapping, categorizing, or structured reporting. If sections shift among different approaches, revise the headings and paragraphs so the manuscript consistently performs the intended type of review work.

\section{Interpretation, Implications, and Limitations}

\subsection{Interpretation}

Interpretation explains what the literature means. It identifies patterns, tensions, gaps, or insights that emerge from the review. Interpretation can appear in a separate discussion section or throughout body sections, depending on the review type and manuscript stage.

\subsection{Implications}

Implications explain why the review matters. They may address:

\begin{itemize}
  \item research: what scholars should study next;
  \item practice: what practitioners might consider;
  \item policy: what decision makers should examine;
  \item theory: how concepts or frameworks may need refinement;
  \item methodology: how future reviews or studies should be designed.
\end{itemize}

\subsection{Limitations}

Limitations explain the boundaries of the manuscript. They should not be presented as failure. They help readers understand scope, method, evidence, and interpretation.

Strong limitations are specific:

\begin{itemize}
  \item The review focuses on peer-reviewed journal articles and does not include dissertations or policy reports.
  \item The search strategy prioritized English-language publications, which may limit international representation.
  \item The manuscript maps emerging themes but does not yet include a full quality appraisal.
\end{itemize}

\subsection{Next Steps}

The final assignment requires a Next Steps component. This section should explain what remains incomplete or developing, what the writer would do next, and how the review could be extended, completed, or strengthened.

\keybox{\textbf{Next Steps should communicate scholarly awareness, not weakness.} A strong Next Steps section shows that the writer understands the current state of the manuscript and the realistic work required to continue developing it.}

\section{Peer Review and the Reviewer Lens}

\subsection{Why Peer Review Matters}

Peer review helps writers see how readers experience the manuscript. A writer may understand why sections belong together, but readers need that logic to be visible on the page. Peer feedback can reveal unclear purpose, weak transitions, underdeveloped analysis, and organizational gaps.

\subsection{Questions a Reviewer Might Ask}

\begin{itemize}
  \item What is the purpose of the manuscript?
  \item Is the selected review methodology clear?
  \item Does the organization fit the methodology?
  \item Where does the manuscript synthesize rather than summarize?
  \item What claims need more evidence?
  \item What terms need clearer definition?
  \item Where do transitions need to explain relationships?
  \item Are limitations and next steps realistic?
\end{itemize}

\subsection{Using Feedback}

Writers are not expected to accept every suggestion. Instead, evaluate feedback by asking:

\begin{itemize}
  \item Does this comment identify a problem that other readers may also experience?
  \item Does the comment connect to the rubric?
  \item Would addressing this issue improve clarity, synthesis, coherence, or alignment?
  \item Is the suggested change compatible with the selected methodology?
\end{itemize}

\section{Audience, Venue, and Publication Fit}

\subsection{Why Venue Fit Matters}

Different journals, conferences, books, and edited collections have different expectations. Fit refers to how well the manuscript aligns with a venue's scope, audience, manuscript type, structure, and style.

Publication expectations can be found in:

\begin{itemize}
  \item aims and scope statements;
  \item author guidelines;
  \item calls for papers;
  \item published articles or chapters;
  \item templates or submission checklists.
\end{itemize}

\subsection{What to Check}

\begin{center}
\begin{tabularx}{0.98\linewidth}{>{\raggedright\arraybackslash}p{0.23\linewidth}Y Y}
\toprule
\textbf{Area} & \textbf{Question} & \textbf{Evidence to inspect} \\
\midrule
Scope and audience & Does my topic and approach belong here? & Aims and scope; recent titles and abstracts. \\
Submission type & Does the venue accept this kind of review? & Author guidelines; calls for papers; published examples. \\
Structure & Does my organization support my method and purpose? & Published articles; journal templates. \\
Formatting and style & What citation, heading, figure, or manuscript format is expected? & Author guidelines; templates. \\
Required components & What must be included before submission? & Submission checklist; author instructions. \\
\bottomrule
\end{tabularx}
\end{center}

\subsection{Positioning the Manuscript}

At this stage, students are not expected to fully prepare a publishable manuscript. The goal is to begin thinking like an author:

\begin{itemize}
  \item Where might this manuscript fit?
  \item Who would read it?
  \item What expectations would the venue have?
  \item What revisions would be necessary for that audience?
\end{itemize}

\section{Targeted Revision Workflow}

\subsection{Activity 1: Manuscript Revision Workflow}

\begin{enumerate}
  \item Diagnose the draft: identify where the purpose is clear or unclear, where synthesis is strong or weak, where connections are missing, whether organization aligns with methodology, and what questions a reviewer might ask.
  \item Select two or three target areas for revision. Include one issue related to organization or coherence, one related to synthesis or analysis, and one related to clarity or scholarly voice.
  \item Revise with purpose by improving transitions, strengthening source connections, and clarifying key claims.
  \item Explain the revisions: what was the issue, what changed, and how the revision improves clarity, synthesis, or coherence.
\end{enumerate}

\subsection{Activity 2: Discussion and Implications}

Revisit the discussion or implications section. Identify one key insight that needs stronger explanation. Add or revise language that explains why the insight matters and how it connects to the broader literature.

\subsection{Activity 3: Limitations and Next Steps}

Revisit the limitations section. Clarify one limitation, explain why it matters, and add one or two realistic next steps. The limitation should connect to the methodology or scope of the review.

\subsection{Activity 4: Publication Positioning}

Identify one possible venue type, describe the likely audience, explain how the manuscript aligns with that audience, and name one revision that would improve fit.

\section{Final Literature Review Manuscript Assignment}

\subsection{Assignment Purpose}

The final manuscript assignment is the culmination of the course manuscript development process. Students submit a revised literature review manuscript that reflects the selected review methodology, demonstrates synthesis and organization, and shows thoughtful revision based on feedback and self-assessment.

Because students are using different review methodologies, the assignment does not require every student to complete every stage of a full publishable review. Instead, the submission should represent the strongest and most developed version of the manuscript at this stage.

\subsection{Required Core Sections}

All manuscripts should include:

\begin{itemize}
  \item \textbf{Introduction:} establishes the topic, purpose, and direction of the review.
  \item \textbf{Methodology or approach:} explains how literature was selected, organized, analyzed, mapped, or interpreted.
  \item \textbf{Organized literature review body:} presents literature through themes, categories, concepts, findings, mapped areas, or another methodology-aligned structure.
  \item \textbf{Reference section:} includes at least 20 scholarly sources, cited and formatted consistently.
\end{itemize}

\subsection{Flexible Developing Functions}

The manuscript should also represent the following functions in ways that fit the selected methodology and current stage:

\begin{itemize}
  \item interpretation of the literature;
  \item implications or contribution;
  \item awareness of limitations or scope;
  \item a required Next Steps component.
\end{itemize}

\subsection{Submission Expectations}

\begin{itemize}
  \item Submit the final manuscript in Word document format.
  \item Include a title or cover page.
  \item Identify and consistently apply an academic formatting style.
  \item Include all planned sections, even if some are still developing.
  \item Include at least 20 scholarly sources.
  \item Expected length: 10 or more double-spaced pages, excluding title page and references.
  \item If using a source manager, submit a version with citations unlinked while keeping a separate working version with active links.
\end{itemize}

\section{Final Manuscript Rubric}

The final manuscript is evaluated on coherence, synthesis, methodology alignment, scholarly voice, and evidence of revision. The total value is 50 points.

\begin{longtable}{@{}>{\raggedright\arraybackslash}p{0.22\linewidth}>{\raggedright\arraybackslash}p{0.40\linewidth}>{\raggedright\arraybackslash}p{0.22\linewidth}>{\raggedright\arraybackslash}p{0.04\linewidth}@{}}
\toprule
\textbf{Criterion} & \textbf{Focus} & \textbf{Strong performance} & \textbf{Pts} \\
\midrule
Core Requirements and Structure & Title/cover page, formatting style, introduction, methodology/approach, literature body, references, minimum 20 sources. & Required components are present and logically structured. & 8 \\
\midrule
Introduction and Purpose & Clear topic introduction, explicit purpose, problem, gap, or focus. & Topic, purpose, and direction are clear and well positioned. & 6 \\
\midrule
Literature Synthesis and Analysis & Patterns, tensions, relationships, gaps, and analysis appropriate to topic and method. & Strong synthesis across sources is clearly identified and connected. & 10 \\
\midrule
Organization and Coherence & Logical flow, headings, subheadings, transitions, and section connections. & Organization is coherent and supports the manuscript purpose. & 6 \\
\midrule
Alignment with Methodology & Structure, analysis, and sections align with the selected review methodology. & Organization and analysis clearly reflect the chosen review type. & 6 \\
\midrule
Interpretations and Implications & Interpretation of findings, mapped areas, implications, or contribution. & Clear interpretation shows emerging or developed understanding of what the literature means. & 5 \\
\midrule
Limitations and Next Steps & Realistic reflection on boundaries, incomplete work, and future development. & Limitations and Next Steps are clear, realistic, and aligned with the manuscript. & 4 \\
\midrule
Use of Evidence and Literature Base & Scholarly sources, supported claims, breadth and depth appropriate to topic and method. & Sources are used strongly and claims are supported. & 3 \\
\midrule
Writing Quality and Scholarly Voice & Clarity, grammar, academic tone, and professionalism. & Writing is clear, professional, and consistent with scholarly expectations. & 2 \\
\bottomrule
\end{longtable}

\keybox{\textbf{Important framing:} A strong final draft is not a perfect manuscript. It is a coherent, methodologically aligned, thoughtfully revised manuscript that demonstrates the ability to synthesize literature and communicate a developing scholarly contribution.}

\section{Revision Checklist}

Use this checklist before submission:

\begin{enumerate}
  \item The introduction states the topic, purpose, and direction.
  \item The methodology or approach is named and explained.
  \item The manuscript body is organized by themes, categories, concepts, findings, mapped areas, or another appropriate structure.
  \item Each section contributes to the manuscript purpose.
  \item Transitions explain relationships between ideas and sections.
  \item Sources are synthesized rather than listed one by one.
  \item Claims are specific, supported, and written in scholarly tone.
  \item Interpretation or discussion explains what the literature means.
  \item Implications or contribution are visible.
  \item Limitations are specific and connected to scope or method.
  \item Next Steps explain how the project could be extended or strengthened.
  \item The reference section includes at least 20 scholarly sources.
  \item Formatting style is consistent.
  \item The manuscript reflects meaningful revision based on feedback and self-assessment.
\end{enumerate}

\section{Key Takeaways}

\begin{itemize}
  \item Week 6 revision focuses on the manuscript as a complete scholarly document.
  \item Coherence depends on purpose, section order, transitions, synthesis, and throughline.
  \item Scholarly voice requires clarity, precision, evidence, and appropriate caution.
  \item Methodology alignment means the manuscript actually performs the work required by the selected review type.
  \item Interpretation, implications, limitations, and next steps help readers understand the value and boundaries of the review.
  \item Peer review should be used strategically to identify high-impact revisions.
  \item The final manuscript is evaluated as a coherent, methodologically aligned, thoughtfully revised literature review in progress.
\end{itemize}

\section{Review Questions}

\begin{enumerate}
  \item What is the difference between proofreading and Week 6 revision?
  \item How can a reverse outline help identify coherence problems?
  \item What is the difference between cohesion and coherence?
  \item How does a synthesized sentence differ from an author-by-author summary?
  \item What does methodology alignment require in a literature review manuscript?
  \item How should a limitations section be framed?
  \item Why is the Next Steps section required in the final manuscript?
  \item What evidence would show that a manuscript fits a possible publication venue?
  \item Which rubric category is most likely to need attention in your current draft?
  \item What two or three revisions would most improve the manuscript before submission?
\end{enumerate}

\section{Condensed Study Sheet}

\begin{center}
\begin{tabularx}{0.98\linewidth}{>{\raggedright\arraybackslash}p{0.24\linewidth}Y}
\toprule
\textbf{Concept} & \textbf{Remember} \\
\midrule
Coherence & The whole manuscript fits together around a clear purpose. \\
Reverse outline & A diagnostic tool that shows what the draft currently does. \\
Scholarly voice & Clear, precise, evidence-based, and appropriately cautious writing. \\
Synthesis & Connecting sources to explain patterns, tensions, relationships, or gaps. \\
Methodology alignment & Purpose, organization, and analysis match the selected review type. \\
Interpretation & Explanation of what the reviewed literature means. \\
Implications & Explanation of why the review matters for research, practice, policy, theory, or method. \\
Limitations & Specific boundaries of the review, framed as scope awareness. \\
Next Steps & Realistic future work needed to extend, complete, or strengthen the manuscript. \\
Reviewer lens & Reading the manuscript as an outside reader who needs the logic made visible. \\
\bottomrule
\end{tabularx}
\end{center}

\end{document}
